\documentclass[11pt,a4paper]{article}

\usepackage[T2A]{fontenc}
\usepackage[utf8]{inputenc}
\usepackage[russian,english]{babel}
\usepackage{geometry}
\usepackage{amsmath,amssymb}
\usepackage{hyperref}
\usepackage{booktabs}
\usepackage{enumitem}

\geometry{margin=2.5cm}

\title{Angel-in-Pocket Full Business Agent with Subscriptions\\
Самоисправляющаяся AI-архитектура бизнеса}
\author{qqewq}
\date{\today}

\begin{document}
\selectlanguage{russian}
\maketitle

\begin{abstract}
Angel-in-Pocket Full Business Agent with Subscriptions — это backend-first AI-система, которая превращает идею пользователя в бизнес, готовый к запуску по модели подписки. Архитектура сочетает мультиагентную оркестрацию, автоматизацию бизнес-процессов и слой устойчивости, вдохновлённый GRA-Core, для оценки сценариев, удаления неэффективных ветвей и сохранения структурной согласованности.
\end{abstract}

\noindent\textbf{Ключевые слова:} ИИ-агенты, автоматизация бизнеса, подписки, устойчивость, оркестрация, финансы, бухгалтерия.

\section{Введение}

Цель Angel-in-Pocket — превратить человеческое видение в работающий бизнес. Пользователь задаёт идею, предпочтения и ограничения, а рой ангельских агентов проектирует продукт, бизнес-модель, подписки и сопутствующие операции.

Система построена вокруг слоя устойчивости, который выявляет противоречия и удаляет неэффективные ветви. Благодаря этому архитектура становится самокорректирующейся, а не просто генеративной.

\section{Постановка задачи}

Многие ИИ-системы способны предлагать идеи, но далеко не все умеют поддерживать целостную бизнес-структуру во времени. Основная задача — преобразовать человеческое видение в прибыльную и устойчивую операционную модель.

Это преобразование можно записать как:
\[
(V, P, C) \rightarrow B \rightarrow G \rightarrow \Pi,
\]
где $V$ — видение, $P$ — модель предпочтений, $C$ — набор ограничений, $B$ — бизнес-канва, $G$ — граф процессов, а $\Pi$ — итоговая операционная политика.

\section{Архитектура}

Репозиторий можно разделить на четыре слоя:

\begin{itemize}[leftmargin=1.2cm]
    \item \textbf{gra\_core}: оценка устойчивости, поиск противоречий, ранжирование сценариев.
    \item \textbf{angel}: агенты для продукта, маркетинга, продаж, финансов, бухгалтерии и рисков.
    \item \textbf{app}: backend-сервисы, API и модели базы данных.
    \item \textbf{frontend}: пользовательский интерфейс и дашборды.
\end{itemize}

Каждый слой участвует в общем цикле управления:
\[
\text{Наблюдение} \rightarrow \text{Оценка} \rightarrow \text{Обнуление} \rightarrow \text{Пересборка}.
\]

\section{Модель подписок}

Бизнес-модель ориентирована на регулярную выручку. Тариф подписки можно представить как:
\[
T_i = \{p_i, b_i, r_i, c_i\},
\]
где $p_i$ — цена, $b_i$ — период списания, $r_i$ — логика удержания, а $c_i$ — риск оттока.

Агент оценивает каждый тариф и оставляет только те, которые удовлетворяют ограничениям прибыльности и устойчивости.

\section{Финансы и бухгалтерия}

Для работы в реальном мире система включает каркас финансового и бухгалтерского учёта. Упрощённая формула прибыли:
\[
\Pi = R - O - \tau,
\]
где $R$ — выручка, $O$ — операционные расходы, а $\tau$ — налоги.

Этот слой можно расширить прогнозом денежных потоков, планом счетов, проводками и логикой, зависящей от налогового режима.

\section{Принцип устойчивости}

Главный принцип архитектуры — обнуление нестабильных ветвей. Процесс должен сохраняться только в том случае, если его полезность превышает издержки:
\[
\text{Сохранить}(x) \iff U(x) - K(x) > 0.
\]

Слой устойчивости сравнивает альтернативные ветви и выбирает ту, которая обладает наибольшей структурной согласованностью.

\section{Сценарии применения}

Система подходит для:

\begin{itemize}[leftmargin=1.2cm]
    \item бизнесов с подписочной моделью,
    \item productized service-систем,
    \item бизнес-планирования и симуляции,
    \item агентной оркестрации с учётом финансов,
    \item автоматизации стартапов с контролем рисков.
\end{itemize}

\section{Заключение}

Angel-in-Pocket — это не просто планировщик или чатбот. Это самоисправляющаяся AI-архитектура бизнеса, в которой человек задаёт видение, ангельские агенты исполняют бизнес-логику, а слой устойчивости сохраняет согласованность системы во времени.

\begin{thebibliography}{9}

\bibitem{latex}
LaTeX Project. \emph{LaTeX — система подготовки документов}. \url{https://www.latex-project.org/}

\bibitem{ctan}
CTAN. \emph{The Comprehensive TeX Archive Network}. \url{https://ctan.org/}

\end{thebibliography}

\end{document}